\documentclass[10pt,a4paper]{amsart}
\usepackage[margin=19mm]{geometry}
\usepackage{amsmath,amssymb,mathtools}
\usepackage[scr=rsfs,cal=euler]{mathalfa}
\usepackage{xfrac}
\usepackage[unicode,hidelinks,bookmarks=false]{hyperref}
\newcommand{\ie}{i.e.\ }
\newcommand{\cf}{cf.\ }
\newcommand{\resp}{respectively\ }
\newcommand{\Subsection}{\subsection}
\providecommand{\nobreakdash}{\nobreak\hbox{-}}
\setlength{\emergencystretch}{2em}
\multlinegap0pt
\usepackage{bm}
\usepackage{bbm}
\usepackage{dsfont}
\usepackage{xspace}
\usepackage{cancel}
\usepackage{mathtools}
\usepackage{amsthm}
\makeatletter
\@ifundefined{theorem}{\newtheorem{theorem}{Theorem}}{}
\@ifundefined{lemma}{\newtheorem{lemma}[theorem]{Lemma}}{}
\@ifundefined{proposition}{\newtheorem{proposition}[theorem]{Proposition}}{}
\makeatother
\DeclareMathOperator{\Res}{Res}
\newcommand{\Dop}{D}
\newcommand{\Hh}{\mathfrak{H}}
\newcommand{\dd}{\mathrm{d}}
\newcommand{\ii}{\mathrm{i}}
\title[Eta-products, Eichler integrals, and the level-8 Apery limit]{Eta-products, Eichler integrals, and the level-8 Apery limit}
\author{Alex Shvets}
\date{Source: arXiv:2604.14219v2 (CC BY 4.0)}
\begin{document}
\maketitle

\noindent\emph{Source and licence.} Eta-products, Eichler integrals, and the level-8 Apery limit. Version arXiv:2604.14219v2 (CC BY 4.0), distributed under the Creative Commons Attribution 4.0 International licence, as recorded in the arXiv licence metadata for this exact version and shown on the abstract page https://arxiv.org/abs/2604.14219v2. Full rights notice: licenses/arxiv-2604.14219-CC-BY-4.0.txt.\par\smallskip

\noindent\textit{Editorial scope.} Counted unit: the Proposition whose source label is \texttt{prop:period-polynomial} in the article (source0.tex lines 627--640, source label \texttt{prop:period-polynomial}), delivered together with the article's own complete original proof of it (source0.tex lines 642--717, one \texttt{proof} environment of 76 lines). No mathematics is written, repaired or re-derived: statement and proof are the article's own text line for line. Required context retained uncounted: prop:E-eichler (lines 451--465); eq:E-qseries (lines 457--460); lem:g-transform (lines 475--480); eq:L-functional (lines 523--526); lem:bol (lines 595--605). Every definition, display and label of those lines is retained unchanged. Omitted from this package and not redistributed: 1--450, 461--474, 481--522, 527--594, 606--626, 641--641, 718--1087. Nothing inside a retained excerpt is omitted or altered: every retained body is delivered exactly as the article's lines are written, so no omission receipt of any kind accompanies this package. Citations: the retained excerpts carry no citation command at all, so no bibliography entry of the article is transcribed and no reference is redistributed. Reference adaptation: none was needed and none was made; every reference inside the delivered unit resolves inside it, so no unresolved reference or citation is produced. Printed slips kept literal: no printed word, symbol, hypothesis or number of the counted statement or of its proof was altered. Layout and class: the article's own class is \texttt{amsart} with packages including fontenc, inputenc, lmodern, amsmath, amssymb, amsthm, hyperref; the wrapper is the lane's pinned \texttt{amsart} editorial wrapper at 10pt on A4, which restates the theorem-like environments the retained text needs and adds the article's own macro definitions that the retained text needs (\texttt{\textbackslash Dop}, \texttt{\textbackslash Hh}, \texttt{\textbackslash Res}, \texttt{\textbackslash dd}, \texttt{\textbackslash ii}), with extra wrapper packages (bm, bbm, dsfont, xspace, cancel, mathtools). The delivered numbering is the unit's own. Assets: no retained span contains a figure environment, a graphics-inclusion command, a PDF-page inclusion, a table or a TikZ picture; no image file is copied into this package, the package declares no asset dependency and no image file is redistributed with it. Licence: version arXiv:2604.14219v2 of this article is distributed under the Creative Commons Attribution 4.0 International licence, as recorded in the arXiv licence metadata for this exact version and shown on the abstract page https://arxiv.org/abs/2604.14219v2. A full rights notice is delivered at licenses/arxiv-2604.14219-CC-BY-4.0.txt. Characters: every retained excerpt is pure ASCII, so the delivered engine, XeLaTeX with its default Latin Modern text font, reports no missing character.
\begin{proposition}\label{prop:E-eichler}
One has
\begin{equation}\label{eq:Phi-equals-g}
\Phi=g_8,
\end{equation}
and therefore
\begin{equation}\label{eq:E-qseries}
E(\tau)=\sum_{n\ge 1}\frac{a_n}{n^3}q^n
=-\frac{(2\pi\ii)^3}{2}\int_{\tau}^{\ii\infty}g_8(z)(z-\tau)^2\,\dd z.
\end{equation}
Equivalently,
\begin{equation}\label{eq:B-param}
\sum_{n\ge 0}B_n^{(8)}t(\tau)^n=E(\tau)Y(\tau).
\end{equation}
\end{proposition}

\begin{lemma}\label{lem:g-transform}
One has
\begin{equation}\label{eq:g-transform}
g_8|_4W_8=-g_8.
\end{equation}
\end{lemma}

\begin{equation}\label{eq:L-functional}
\Gamma(s)\left(\frac{\pi}{\sqrt2}\right)^{-s}L(g_8,s+3)
=\Gamma(-s-2)\left(\frac{\pi}{\sqrt2}\right)^{s+2}L(g_8,1-s).
\end{equation}

\begin{lemma}\label{lem:bol}
For every holomorphic function $F$ on $\Hh$ one has
\begin{equation}\label{eq:bol}
\frac{\dd^3}{\dd\tau^3}\bigl(8\tau^2F(W_8\tau)\bigr)
=\frac{1}{64\tau^4}F'''(W_8\tau).
\end{equation}
Equivalently,
\[
\Dop^3(F|_{-2}W_8)=(\Dop^3F)|_4W_8.
\]
\end{lemma}

\begin{proposition}[Explicit period polynomial]\label{prop:period-polynomial}
Let
\begin{equation}\label{eq:P8-def}
P_8(\tau):=(E|_{-2}W_8)(\tau)+E(\tau).
\end{equation}
Then $P_8$ is a polynomial of degree at most $2$, and in fact
\begin{equation}\label{eq:period-polynomial}
P_8(\tau)=\frac{7}{32}\zeta(3)(8\tau^2+1).
\end{equation}
Equivalently,
\begin{equation}\label{eq:E-transform}
(E|_{-2}W_8)(\tau)+E(\tau)=\frac{7}{32}\zeta(3)(8\tau^2+1).
\end{equation}
\end{proposition}

\begin{proof}
Set
\[
H(\tau):=(E|_{-2}W_8)(\tau)+E(\tau).
\]
By Lemma~\ref{lem:bol}, Proposition~\ref{prop:E-eichler}, and Lemma~\ref{lem:g-transform},
\[
\Dop^3H=(\Dop^3E)|_4W_8+\Dop^3E=g_8|_4W_8+g_8=0.
\]
Hence $H$ is a polynomial of degree at most $2$.

To determine it explicitly, restrict to the $W_8$-invariant geodesic
\begin{equation}\label{eq:tauy-def}
\tau(y):=\frac{\ii y}{2\sqrt2}
\qquad(y>0).
\end{equation}
Then $W_8\tau(y)=\tau(1/y)$, and \eqref{eq:E-qseries} gives
\begin{equation}\label{eq:Fy-def}
F(y):=E(\tau(y))=\sum_{n\ge 1}\frac{a_n}{n^3}\exp\!\left(-\frac{\pi ny}{\sqrt2}\right).
\end{equation}
Fix $1<c<2$. Mellin inversion yields
\begin{equation}\label{eq:mellin}
F(y)=\frac{1}{2\pi\ii}\int_{(c)}
\Gamma(s)\left(\frac{\pi}{\sqrt2}\right)^{-s}L(g_8,s+3)y^{-s}\,\dd s,
\end{equation}
because
\[
\int_0^{\infty}\exp\!\left(-\frac{\pi ny}{\sqrt2}\right)y^{s-1}\,\dd y
=\Gamma(s)\left(\frac{\pi n}{\sqrt2}\right)^{-s}
\qquad(\Re(s)>0).
\]
By Stirling's formula, the integrand in \eqref{eq:mellin} decays exponentially on vertical lines. We may therefore shift the contour to $\Re(s)=-c-2$. In the strip $-c-2<\Re(s)<c$, the integrand $M(s)y^{-s}$ has simple poles only at $s=0,-1,-2$ (from $\Gamma(s)$); the pole of $\Gamma$ at $s=-3$ is cancelled by the zero $L(g_8,0)=0$. The residue at $s=-1$ vanishes because $L(g_8,2)=0$.

Write
\[
M(s):=\Gamma(s)\left(\frac{\pi}{\sqrt2}\right)^{-s}L(g_8,s+3).
\]
Using the functional equation \eqref{eq:L-functional} and the substitution $u=-s-2$, we obtain
\begin{align*}
\frac{1}{2\pi\ii}\int_{(-c-2)}M(s)y^{-s}\,\dd s
&=\frac{1}{2\pi\ii}\int_{(c)}M(u)y^{u+2}\,\dd u\\
&=y^2F(1/y).
\end{align*}
Hence the residue theorem gives
\begin{equation}\label{eq:residue-identity}
F(y)-y^2F(1/y)
=\Res_{s=0}\bigl(M(s)y^{-s}\bigr)+\Res_{s=-2}\bigl(M(s)y^{-s}\bigr).
\end{equation}
The residues are
\[
\Res_{s=0}\bigl(M(s)y^{-s}\bigr)=L(g_8,3)=\frac{7}{32}\zeta(3),
\]
\[
\Res_{s=-2}\bigl(M(s)y^{-s}\bigr)
=\frac12\left(\frac{\pi}{\sqrt2}\right)^2L(g_8,1)y^2
=-\frac{7}{32}\zeta(3)y^2.
\]
Substituting into \eqref{eq:residue-identity} yields
\begin{equation}\label{eq:F-functional}
F(y)-y^2F(1/y)=\frac{7}{32}\zeta(3)(1-y^2).
\end{equation}
On the geodesic \eqref{eq:tauy-def},
\[
(E|_{-2}W_8)(\tau(y))=8\tau(y)^2E(W_8\tau(y))=-y^2F(1/y),
\]
so \eqref{eq:F-functional} becomes
\[
H(\tau(y))=\frac{7}{32}\zeta(3)(1-y^2).
\]
Since $8\tau(y)^2+1=1-y^2$, this is exactly
\[
H(\tau(y))=\frac{7}{32}\zeta(3)(8\tau(y)^2+1)
\qquad(y>0).
\]
Both sides are polynomials of degree at most $2$ in $\tau$ and they agree on the infinite set $\{\tau(y):y>0\}$, so they are identical. This proves \eqref{eq:period-polynomial} and \eqref{eq:E-transform}.
\end{proof}

\end{document}
